% Generated draft: requires mathematical and source review.
\documentclass[11pt]{article}
\usepackage[margin=25mm]{geometry}
\usepackage{fontspec}
\setmainfont{Latin Modern Roman}
\usepackage{amsmath,amssymb,mathtools,unicode-math}
\setmathfont{Latin Modern Math}
\usepackage{enumitem,needspace,xurl,hyperref}
\setlength{\parindent}{0pt}
\setlength{\parskip}{5pt}
\setlength{\emergencystretch}{3em}
\newcommand{\N}{\mathbb{N}}
\newcommand{\Z}{\mathbb{Z}}
\newcommand{\Q}{\mathbb{Q}}
\newcommand{\R}{\mathbb{R}}
\newcommand{\C}{\mathbb{C}}
\newcommand{\F}{\mathbb{F}}
\newcommand{\A}{\mathbb{A}}
\newcommand{\PP}{\mathbb{P}}
\newcommand{\E}{\mathbb{E}}
\newcommand{\eps}{\varepsilon}
\newcommand{\dd}{\,\mathrm{d}}
\begin{document}
\section*{UnsolvedMath ID: 30006328}
\textbf{Source problem code:} catalogue entry
\subsection{Definitions and mathematical statement}
The supplied record identifies the intended assertion as the \emph{Geometric Numerical Conjecture}, cited in the excerpt as [2, 14.31]. Its formulation is conditional on preceding notation and hypotheses that are absent from the supplied extract. Accordingly, the following records the recoverable assertion and explicitly marks the unresolved right-hand side.

Write \(\F_q\) for the finite field of cardinality \(q\). Let \(\mathcal F_\phi\) be a cuspidal Hecke eigensheaf in the setting intended by the source, with trace function \(f_\phi\). The notation \(\operatorname{Bun}_G\) denotes the moduli stack of \(G\)-bundles on the source's unspecified base curve. Write \(P_X^{*,\mathrm{norm}}\) for the trace function of the normalized \(*\)-period sheaf. The excerpt describes the \(*\)-period sheaf as the naive Verdier dual of the \(!\)-period sheaf and warns that this trace function differs from \(P_X^{\mathrm{norm}}\) in its Conjecture 2; it refers to [2, 14.8] for the distinction.

The intended assertion is that, for every such cuspidal Hecke eigensheaf, its normalized global period satisfies an identity of the form
\[
 \sum_{x\in\operatorname{Bun}_G(\F_q)}
 f_\phi(x)P_X^{*,\mathrm{norm}}(x)
 \;\stackrel{?}{=}\;
 q^{-b_G/2}\,\mathcal S_{\mathrm{source}}(\phi^d).
\]
Here \(\mathcal S_{\mathrm{source}}(\phi^d)\) is an explicit placeholder for the unrecoverable normalized sheared \(L\)-expression in the source, evaluated at \(s=1\); it is not a newly defined invariant. The corrupted extraction contains \(T^\vee\), \(\check X\), and an apparent index \(x\in\check X_\phi\), but does not reliably determine the coefficient object, its evaluation point, or whether an additional sum is required. Neither \(b_G\) nor the precise normalization or summation convention on the left is supplied. Thus this display records the intended identity's shape, not a complete mathematical formulation.

The excerpt describes \(d\), following its reference [8], as an involution associated with a pinning of \(G\): it negates the pinning on the simple roots and acts on the torus by \(t\mapsto w_G(t^{-1})\), where \(w_G\) is a lift of the longest Weyl element taking \(B\) to \(B^-\). For \(G=\mathrm{SL}_n\), it gives
\[
 g\longmapsto\operatorname{Ad}(w_G)\bigl({}^{\mathsf t}(g^{-1})\bigr).
\]
The notation \(\phi^d\) denotes the resulting twist of the Galois parameter: the excerpt describes the action of \(\Gamma_F\) on \(\check X\), through \(\check G\times G_{\mathrm{gr}}\), as precomposed with \((-)^d\). The precise group carrying this involution and the requisite identifications are not supplied.

For a \(\Gamma_F\times\mathbb G_m\)-module \(V\) under the action specified by \(\phi\), the excerpt gives a grading
\[
 V=\bigoplus_{i\in\Z}V^{(i)},
\]
where each \(V^{(i)}\) is a \(\Gamma_F\)-submodule of grading weight \(i\). Its displayed prescription for the normalized sheared \(L\)-function is
\[
 L^{\mathrm{norm}}(s,\phi,V)
 =\prod_i L^{\mathrm{norm}}(s,V^{(i)}),
\]
with the factors described as Artin \(L\)-functions normalized according to [2, Chapter 11]. The shearing notation is corrupted in the extraction; no grading-dependent shifts, normalization conventions, or conditions making this product meaningful can be recovered reliably.

The supplied curation describes the trace-of-Hom reduction, citing [2, Lemma 2.6.1], as formal, while identifying the resulting arithmetic identity as conjectural in general. This draft does not establish either claim and requires the missing source context before a fully quantified statement can be given.
\subsection{Short English statement}
The intended conjecture equates the normalized global period of every eligible cuspidal Hecke eigensheaf with a normalized sheared \(L\)-expression for its duality-twisted parameter, including the factor \(q^{-b_G/2}\). The supplied extraction omits essential hypotheses and corrupts the right-hand side, so the precise identity remains unresolved. The supplied open status and description of a formal trace-of-Hom reduction are curation claims, not independently verified conclusions.
\subsection{Sources}
\begin{itemize}
\item Supplied local UnsolvedMath registry record, permanent ID \(30006328\), title \emph{Arithmetic Identities from Geometric Trace Formulas}, source problem code \texttt{catalogue entry}, in the Oberwolfach Reports Open Problems collection. No individual proposer is supplied; the proposed year is 2025.
\item \emph{Arbeitsgemeinschaft: Relative Langlands Duality} (2025), Oberwolfach Reports, pp. 811--882, DOI \texttt{10.4171/owr/2025/18}. Supplied report citation and literature-source link: \url{https://doi.org/10.4171/owr/2025/18}.
\item \emph{Relative Langlands Duality}. Supplied literature-source link: \url{https://arxiv.org/abs/2409.04677}. Authors and a precise correspondence with the excerpt's numbered bibliography are not supplied.
\item The excerpt attributes the conjecture to [2, 14.31], the trace-of-Hom calculation to [2, Lemma 2.6.1], the period-sheaf distinction to [2, 14.8], the Artin \(L\)-function normalization to [2, Chapter 11], and the duality involution to [8]. These reference labels are retained; their complete bibliography is absent.
\item Literature assessment checked at \texttt{2026-08-22}: the supplied curation labels the problem open, describes the trace-of-Hom reduction as formal, and states that the resulting arithmetic identity is not proved in general. This is supplied dated curation, not independently verified. The record also labels its statement \texttt{corrected\_verified} and reports that the official report was checked; those are supplied verification claims, not verification performed for this draft.
\end{itemize}
\end{document}
